\subsection*{49. Work Samples Ask Candidates to Do Recognizable Parts of the Job}

A \textbf{work sample} requires performance of a task substantially resembling actual work. A writer edits a flawed memo. An analyst examines a dirty dataset and explains uncertainty. An instructor diagnoses a misconception and gives feedback. An engineer debugs a realistic failure. A service manager prioritizes competing customer incidents. This directness makes work samples attractive, and applicants often regard them as intuitively relevant when designed fairly [R7].

But a work sample is not automatically valid. It may be too narrow, easy to game, dependent on an unfamiliar tool, biased toward candidates who have encountered the exact problem before, or unlike the job's normal environment. Its task content must come from job analysis; its scoring must be consistent; and its burden should be proportionate.

Work samples are most appropriate for capabilities expected upon entry. If the firm plans to train everyone in a specialized internal workflow, testing detailed pre-existing familiarity may simply reward outsiders who happen to know the firm's chosen tool, rather than future success.

\subsection*{50. Fidelity Requires Representative Constraints, Not Theatrical Realism}

A realistic simulation includes the features of work that change the decisions being measured. It need not reproduce every physical detail. An engineer often uses documentation, code search, tests, and AI assistants; banning all tools may assess memory rather than engineering judgment. A reviewer usually receives context and can ask clarifying questions; withholding all context may measure tolerance for ambiguity rather than review quality.

Conversely, a candidate who receives unlimited time on a take-home exercise may produce a beautifully polished artifact that conceals how they prioritize under normal delivery constraints. Choose time and tool rules that resemble the actual job. Publish them consistently. Include 1 part of the assessment where the candidate must explain decisions, demonstrate understanding, or respond to a realistic change, especially when generated output is easy to obtain.

The strongest work sample is \textbf{representative enough to reveal the required judgment, short enough to respect candidates, and controlled enough to permit comparison}. Those goals sometimes conflict; record the compromise instead of pretending the test mirrors every aspect of employment.

\subsection*{51. Sample the Work Broadly Enough to Avoid Overfitting}

Suppose an instructor performs a brilliant 10-minute lecture on a familiar topic. That is evidence of preparation and delivery. It is not evidence that the instructor can diagnose a confused student, give feedback on flawed work, keep a mixed-ability group engaged, or assess honestly. A single polished performance can be a narrow measurement of a broad role.

1 alternative is a small portfolio of tasks: explain a concept, respond to an unscripted misconception, mark an authentic-looking submission, and plan how to remediate a struggling learner. The tasks together sample more of the teaching process. For an engineer, the corresponding set might include specification, debugging, evaluation, and a design defense.

Task sampling improves coverage but has a cost. The founder should identify \textbf{the minimum representative set} that changes a decision. If every competent applicant excels at a particular exercise, it may have little power to discriminate among finalists. Revise it based on observed outcomes rather than continuously adding more tasks.

\subsection*{52. Short Take-Home Tasks Can Be Fair; Unpaid Production Work Is Not}

An assessment that asks candidates to complete 4 hours of research, design a full syllabus, produce marketing assets, and solve the company's current operational problem risks becoming uncompensated consulting. Even if the founder has no intention of using the results commercially, substantial unpaid work imposes a burden that candidates bear unequally.

Use self-contained, fictional, or de-identified tasks designed for evaluation. Short tasks may reasonably be unpaid where lawful, customary, and genuinely proportionate. Lengthier finalists' exercises should often be compensated. A paid trial must have clear terms about deliverables, compensation, ownership, confidentiality, hours, and whether any actual customer work is involved. A candidate should not need to disclose confidential former-employer documents to demonstrate capability.

A small startup tempted to economize by extracting free samples is especially short-sighted. The best candidates can decline such arrangements, and the institution's future reputation may suffer. \textbf{The cost of assessment belongs partly to the employer that benefits from the decision.}

\subsection*{53. Dynamic Work Samples Can Reveal How People Learn}

Some jobs involve unfamiliar tools, changing requirements, and incomplete information. A static test of current knowledge may underpredict candidates who learn quickly. A \textbf{dynamic assessment} provides an initial task, targeted feedback or new information, and a related 2nd task. Evaluators observe whether the candidate understands the correction, revises the approach, and transfers the principle rather than copying the previous solution.

For example, an AI engineering candidate may initially select an inappropriate evaluation metric. The assessor explains a distribution-shift problem. The candidate must design a modified evaluation and justify what changed. The aim is to observe learning and judgment, not to humiliate someone for an initial error. All candidates should receive equivalent information or coaching; otherwise the intervention becomes a hidden source of unfairness.

Such assessments are promising for learning-intensive roles but should not be described as possessing proven universal predictive superiority merely because they seem sophisticated. Treat their predictive utility as a hypothesis to validate. They are especially useful when the company is willing and able to train capability after entry.

\subsection*{54. Assess Critical Ethical and Safety Decisions Without Moral Theater}

Many roles contain decisions in which technical competence and integrity interact: whether to conceal a failed experiment, manipulate a performance report, overstate a learner's results, share confidential information, or promise a customer an impossible outcome. These are legitimate assessment targets when grounded in actual work.

A good scenario asks what the candidate would do, which evidence they would preserve, whom they would notify, and how they would balance competing obligations. A weak test asks broad declarations such as \textit{Are you ethical?} Nearly everyone knows the desired answer. A candidate who identifies uncertainty and a sensible escalation route may be stronger than 1 who promises heroic individual certainty.

Do not require disclosure of sensitive private beliefs or personal life as a proxy for values. The organization should assess \textbf{job-relevant professional conduct}. It must also make its own expectations consistent: screening for honest escalation while penalizing employees who report bad news creates a culture whose incentives contradict its selection claims.
